
\begin{center}
    {\Large \textbf{FUNCTIONS OF SEVERAL VARIABLES}}\\[2mm]
(DEFINITION, DOMAIN, RANGE, GRAPHS AND LEVEL CURVES) \\[4mm]

\mycenterbox{5 minute review}%\\[4mm]
Recap the definition of a function of several variables, and the meaning of domain, range, graph and level curves of functions of several variables.
\end{center}
\mycenterbox{Warm-up.}%\\[4mm]
\problemAnswer{
Match the function with its graph (labeled I–VI).Give reasons for your choices
\begin{multicols}{3}
\begin{enumerate}[label=(\alph*)]
\item $f(x,y)=|x|+|y|$\\[4mm]
\item $f(x,y)=|xy|$
\item $f(x,y)=\frac{1}{1+x^2+y^2}$\\2mm]
\item $f(x,y)=(x^2-y^2)^2$
\item $f(x,y)=(x-y)^2$\\[4mm]
\item $f(x,y)=\sin(|x|+|y|)$
\end{enumerate}
\end{multicols}

\begin{center}
%\centering
\includegraphics[width=10cm]{MCTutorial1Q1.JPEG}
\end{center}
%\tcbox[tcbox raise base]{Warm-up}
 
%\begin{flushright}
%\begin{tikzpicture}
%    % A body
%
%
%    % Non-rectangle shapes must be drawn as polygone
%    \begin{scope}[shift={(3,0)}]
%        \draw  (0,0.5) -- (0,4) -- (.5,4) -- (.5,4.5) -- (1.,4.5) -- (4.,4.5) --(4.,4)-- (4.5,4) --(4.5,0.5) --(4.0,0.5)--(4.0,0.0)-- (.5,0) --(0.5,0.5) -- cycle;
%        % Draw symmetry lines
%        \draw [symmetry] (.5, 0.5) -- (.5,4.0);
%        \draw [symmetry] (4, 0.5) -- (4,4.0);
%        \draw [symmetry] (.5,4.0) -- (4,4.0);
%         \draw [symmetry] (.5, 0.5) -- (4, 0.5);
%        % Dimensions
%        \draw (4.5,0) -- ++(0.5,0) coordinate (D1) -- +(5pt,0);
%        \draw (4.5,4.5) -- ++(0.5,0) coordinate (D2) -- +(5pt,0);
%        \draw [dimen] (D1) -- (D2) node {6cm};
%       
%        
%        \draw (0.0,5) -- ++(0,0.5) coordinate (D1) -- +(0,+5pt);
%        \draw (0.5,5) -- ++(0,0.5) coordinate (D2) -- +(0,+5pt);
%        \draw [dimen,-] (D1) -- (D2) node [above=5pt] {$x$ cm};
%        \draw [dimen,<-] (D1) -- ++(-5pt,0);
%        \draw [dimen,<-] (D2) -- ++(+5pt,0);
%    \end{scope}
%\end{tikzpicture}
%\end{flushright}

}


\mycenterbox{Problems.}%

%\begin{center}
%\tcbox[tcbox raise base]{ Problems. Choose from the below}
%\end{center}
\problem{Exploring functions}%
\begin{enumerate}[(a)]
\item Are the following functions even, odd, or neither? What can you say
about their range?
\begin{multicols}{4}
 \begin{enumerate}[(i)]
\item $\displaystyle \frac{3x}{x^2+2}$
\item $\displaystyle \cos x+3$
\item $\displaystyle \frac{2x^3}{\cos x+3}$
\item $\displaystyle \frac{3x^4+2x^2}{\sin x+4}$
\end{enumerate}
\end{multicols}
\item Consider the functions $\sin^n x$ and $\cos^n x$, where $n$ is a positive integer.
Which are odd, which are even, and what are their periodicities?
\item What is the periodicity of $f(x) = 2 \cos x+\sin 2x$? Is it even/odd/neither?

\end{enumerate}

\problem{Cubic functions} 
Show that $x-1$ is a factor of the cubic function%
$$P(x)=x^3-2x^2-5x+6,$$
and  find the quadratic factor $Q(x)$ such that $P(x) = (x -1)Q(x)$. Factorise
$Q(x)$ and hence sketch $P(x).$

%\newpage
\problem{Sigmoid functions*.}%
For positive integers $n$, consider the family of functions
$$\displaystyle f_n= \frac{x^n}{2^n+x^n}.$$
\begin{enumerate}[(a)]
\item What are appropriate domains and ranges for $f_n(x)$? Do they depend on $n$?

\item Is $f_n(x)$ even, odd, or neither? Does this depend on $n$?

\item Show that for $x\geq 0$, $f_n(x)$ are strictly increasing functions and that
$0\leq f_n(x) < 1.$
\item What is the value of $f_n(2)$? Does it depend on $n$?

\item By considering the values of $f_n(1.9)$ and $f_n(2.1)$ for $n = 2, 4, 8, 16, 32$ and
$64$, show that the graph of $f_n(x)$ for $x\ge 0$ is an increasing sigmoid (that
is, $S$-shaped) function that approaches a step function as $n\rightarrow \infty$.
\end{enumerate}

\mycenterbox{Selected answers and hints}%

For the warm-up, $V(x)=4x(3-x)^2$. Writing $x=1+\epsilon$, we get
$$V = 4(1 +\epsilon)(3- (1+\epsilon))^2=4(4-\epsilon^2(3-\epsilon)).$$
If $-1<\epsilon<2$ then $\epsilon^2 (3-\epsilon)\ge 0$, with equality when $\epsilon=0$. Thus the maximum
value of $V$ occurs when $\epsilon= 0$, i.e. $x = 1\rm{cm}$.

\begin{enumerate}
\item 
\begin{enumerate}[(i)]
\item Odd. The range is the interval  $[-\frac{3}{4}\sqrt{2},\; \frac{3}{4}\sqrt{2}]$
\item Even. The range is the interval $[2, 4]$.
\item Odd. The range is all of $\mathbb R$.
\item Neither. (For example, writing $\displaystyle \frac{5}{\sin x+4}$ we have $\displaystyle f(1)= \frac{5}{\sin (1)+4} \equiv 1.032 (3\rm{dp})$ which is neither $f(1)$ nor $-f(1)$.)\\[4mm]
The range is the nonnegative reals $[0,\infty)$

\end{enumerate}
\end{enumerate}

